\documentclass[12pt,a4paper]{article}
\usepackage[T1]{fontenc}\usepackage{lmodern}
\usepackage[margin=22mm]{geometry}
\usepackage{amsmath,bm,xcolor,titlesec,fancyhdr}
\definecolor{ink}{RGB}{31,42,68}\definecolor{accent}{RGB}{139,30,45}\definecolor{draft}{RGB}{22,86,140}
\setlength{\parindent}{0pt}\setlength{\parskip}{7pt}\linespread{1.25}
\pagestyle{fancy}\fancyhf{}\renewcommand{\headrulewidth}{0pt}
\fancyfoot[L]{\small\color{ink!70}Thesis predefense: speaker notes}\fancyfoot[R]{\small\color{ink!70}\thepage}
\newcommand{\cue}[1]{{\itshape\color{accent}[#1]}}
\newcommand{\drafted}[1]{{\color{draft}#1}}
\newcommand{\slide}[3]{\par\vspace{10pt}\noindent{\large\bfseries\color{ink}Slide #1\quad #2}\hfill{\small\color{ink!70}#3}\par
  \vspace{-4pt}{\color{accent}\rule{\linewidth}{0.6pt}}\par}
\begin{document}\color{ink}
{\LARGE\bfseries Speaker notes}\par\vspace{2pt}
{\color{ink!75}Modeling and Control of Unsteady Flows over Rapidly Maneuvering Airfoils at Ultra-Low Reynolds Number}\par
\vspace{6pt}{\small \cue{Italic in brackets}: pointing cues.\quad \drafted{Blue text}: drafted by Claude where the notes were empty.}\par
\slide{1}{Cover}{}
\drafted{Good morning. I am Imdadullah Raji, and together with Abdullah Al Mamun, under the supervision of Dr Md Ali, I will present our thesis: Modeling and Control of Unsteady Flows over Rapidly Maneuvering Airfoils at Ultra-Low Reynolds Number.}

\slide{2}{Opening quote}{}
The Navier--Stokes equations give us the highest-fidelity representation of fluid flows. However, they give us surprisingly little insight into the flow physics in question. As Feynman asks, ``could we predict the auroras?''

The question that concerns us here is: could we see the vortices roll up like that just by looking at the equations? No, but we can indeed calculate them, like this animation, which came out of a calculation that took 2 hours on 6 cores.

\slide{3}{Objective: Reduced-Order Models}{}
However, we can only afford fractions of a millisecond when it comes to active flow control. Therefore, the goal is to obtain a simplified model. One way to do this is to resolve the flow into patterns observed in space. Proper orthogonal decomposition, shown here, is one such technique. It describes the flow as a sum of its most energetic modes. The simulation data takes gigabytes on our computers, but this POD takes only about 12 megabytes, and a fraction of a millisecond to compute: 19 $\mu$s on a single core of my machine, to be exact.

\slide{4}{Objective: Model-Based Active Flow Control}{}
After reducing the flow to a smaller state vector x, we need a model for how those variables evolve under actuation: the f. \cue{Point to the board.}

We first investigate whether a simple linear model is sufficient. \cue{Point to the linear equation.}

As the dynamics are in fact nonlinear, we can expect that to fail. Then nonlinear system-identification tools can be employed.

For example, Sparse Identification of Nonlinear Dynamics, or SINDy for short, identifies a nonlinear model directly from simulation data. The ultimate purpose is to obtain a controller that maps the state x to an actuation u. \cue{Again point to the relevant equation.}

\slide{5}{Numerical Solver}{20 s}
The simulations were conducted in the open-source finite-volume code OpenFOAM. The problem concerns unsteady fluid flows at low Reynolds number. The flow is laminar, hence the smallest scale to resolve is the boundary-layer thickness $\delta$. We solve the Navier--Stokes equations directly, without any turbulence model. The flow is also assumed to be incompressible. Second-order schemes were used for both convection and time stepping.

\slide{6}{Computational Mesh}{25 s}
The simulation domain is a disk with a radius of 30 chord lengths. When the angle of attack needs changing, we rotate the freestream instead of the airfoil. The first-layer height is taken to be approximately one hundredth of the Blasius flat-plate boundary-layer thickness. Together with the growth rate, it ensures that the boundary layer is resolved by approximately 40 cells.

\slide{7}{Pitching Motion: ALE Formulation}{}
The pitch motion is handled by rotating the entire mesh relative to the freestream. This introduces an additional convective term in every cell, due to the mesh motion. \cue{Point to $\mathbf{u}_m$.} This numerical technique is called the Arbitrary Lagrangian--Eulerian method, as it solves the Eulerian Navier--Stokes equations with an arbitrary Lagrangian motion of the cells.

\slide{8}{Wake Patterns at Ultra-Low Reynolds Number}{}
The violent mixing of turbulent flow is absent in laminar flows at low Reynolds number. As a result, large vortices can survive for a long time without breaking up into smaller ones. In this regime, the shear layer is thick, the wake is broad, and large vortices interact with each other to produce such intricate flow structures. This structure was studied computationally by Kurtulu\c{s}. She studied a NACA0012 airfoil and classified the observed flow structures into five modes according to the wake vorticity pattern.

First, the attached flow. Second, a von K\'arm\'an-like vortex street. Third, the alternating vortex-pair shedding mode. Fourth, alternate vortex with single vortex shedding. Fifth, bluff-body vortex shedding.

\drafted{Our own simulations at Re~=~500 reproduce mode 2 at $\alpha$~=~14$^\circ$ and mode 3 at $\alpha$~=~26$^\circ$.} \cue{Point to the two videos.}

\slide{9}{Period Doubling, Tripling and Chaos}{}
Looking at the force histories also reveals the temporal pattern. Durante et al.\ looked at the time series of the force coefficients at Re~=~1000, their phase portraits and their Fourier amplitudes.

Initially the force peaks are single, and the phase portrait is a single loop. Then, near $\alpha$~=~22$^\circ$, the signal transitions into alternating small and large peaks, and the phase portrait contains two loops. The Fourier spectrum develops an amplitude at the subharmonic.

Additionally, period-three and aperiodic, chaotic flows are also observed.

\slide{10}{Beyond 2D: Spanwise Instability}{}
In reality, a 3D spanwise instability develops, in the same way that the symmetric wake of a cylinder goes unsteady and develops the von K\'arm\'an vortex street.

This instability develops as corrugations of different wavelengths in the spanwise direction. The photo here is from Williamson's 1996 review paper on the vortex dynamics of the cylinder wake. \cue{Point to the photos.}

Gupta et al.\ investigated the 3D instability at Re~=~1000 through numerical studies and water-channel experiments. They observed instability mechanisms similar to those of cylinder wakes. They found that the onset of the 3D instability varies as the inverse square root of the Reynolds number, which sits before all the period-doubling bifurcations reported in 2D simulations.

\slide{11}{Pitching Airfoil and the Leading-Edge Vortex}{}
\drafted{When the airfoil pitches up rapidly, the lift has two sources. The first is added-mass, or non-circulatory, lift: the airfoil has to accelerate the fluid around it, so this part appears only while the pitch rate and acceleration are non-zero, the spikes at the start and end of the ramp in the video.} \cue{Point to the added-mass equation and the $\dot\alpha$, $\ddot\alpha$ graphs.}

\drafted{The second is circulatory lift, set by the circulation through the Kutta--Joukowski theorem, $L'$~=~$\rho$U$\Gamma$. During the ramp a leading-edge vortex rolls up and stays over the suction side, adding circulation and therefore lift. Once the motion stops, the vortex detaches and convects downstream, and the lift falls.}

\drafted{Eldredge and Jones show the effect of pitch rate on a flat plate: the faster the manoeuvre, the tighter and stronger the leading-edge vortex, and the higher the lift peak, from about 3 at K~=~0.1 to about 16 at K~=~1.} \cue{Point to the figure.}

\slide{12}{Books Consulted and Code}{}
\drafted{These are the books we consulted. The code developed for this work is available in these two repositories.}

\slide{13}{Questions?}{}
\drafted{Thank you for your attention. We are happy to take your questions.}
\end{document}
